\section{The Art of Controversy}

\begin{quotex}
What is required is not the courage of one's convictions, but rather the courage for an attack upon one's convictions.
\flright{\textsc{Friedrich Nietzsche}}
\end{quotex}


Although discussions of deep and compicated issues are difficult in person, they are especially problematic on the Internet. There is the tendency to excessive polemics, which manifests as the prepensity to defend a particular opinion rather than to be challenged to reconsider one's own opinion. This often results in the magnification of rather small differences.

As Evola points out in \textit{The Individual and the Becoming of the World}\footnote{\url{https://www.gornahoor.net/individual/index.htm}}, growth in consciousness requires a self-emptying of all our opinions, whether from everyday life, religion, or science.

\begin{quotex}
If anyone in a discussion with us is not concerned with adjusting himself to truth, if he has no wish to find the truth, he is intellectually a barbarian. That, in fact, is the position of the mass-man when he speaks, lectures or writes.
\flright{\textsc{José Ortega y Gasset}}
\end{quotex}


Clearly the intent to find the truth must be there. Then it is possible to advance. The mass-man seems to think that is sufficient to say ``I believe this'' or ``I think that'' without any warrant. Oddly enough, this is usually the position of those who ostensibly reject ``egalitarianism'' as such, yet remain staunchly egalitarian in the intellectual realm.

In the following selection from \emph{The Art of Controversy}, Schopenhauer proposes some guidelines for discussion. A small group of philosophical friends can engage in worthwhile discussions, with those ahead leading those behind. \textbf{This would be the seed for the formation of an \emph{elite} given the right circumstances and right motivation.} For the rest, I wish you peace, the peace to go to sleep every night secure in your beliefs.

\begin{quotationx}
As a sharpening of wits, controversy is often, indeed, of mutual advantage, in order to correct one's thoughts and awaken new views. But in learning and in mental power both disputants must be tolerably equal: If one of them lacks learning, he will fail to understand the other, as he is not on the same level with his antagonist. If he lacks mental power, he will be embittered, and led into dishonest tricks, and end by being rude.

The only safe rule, therefore, is that which Aristotle mentions in the last chapter of his \emph{Topica}: not to dispute with the first person you meet, but only with those of your acquaintance of whom you know that they possess sufficient intelligence and self-respect not to advance absurdities; to appeal to reason and not to authority, and to listen to reason and yield to it; and, finally, to cherish truth, to be willing to accept reason even from an opponent, and to be just enough to bear being proved to be in the wrong, should truth lie with him. From this it follows that scarcely one man in a hundred is worth your disputing with him. You may let the remainder say what they please, for every one is at liberty to be a fool --- \emph{desipere est jus gentium}. Remember what Voltaire says: \emph{La paix vaut encore mieux que la vérité} ``Peace is worth more than the truth'' . Remember also an Arabian proverb which tells us that on the tree of silence there hangs its fruit, which is peace.
\flright{\textsc{Arthur Schopenhauer}}
\end{quotationx}



\flrightit{Posted on 2009-08-09 by Cologero}

\begin{center}* * *\end{center}

\begin{footnotesize}
\begin{sffamily}
\texttt{x on 2012-04-08 at 15:16 said:}

``If anyone in a discussion with us is not concerned with adjusting himself to truth, if he has no wish to find the truth, he is intellectually a barbarian. That, in fact, is the position of the mass-man when he speaks, lectures or writes.''

``As a sharpening of wits, controversy is often, indeed, of mutual advantage, in order to correct one's thoughts and awaken new views. But in learning and in mental power both disputants must be tolerably equal: If one of them lacks learning, he will fail to understand the other, as he is not on the same level with his antagonist. If he lacks mental power, he will be embittered, and led into dishonest tricks, and end by being rude.

The only safe rule, therefore, is that which Aristotle mentions in the last chapter of his Topica: not to dispute with the first person you meet, but only with those of your acquaintance of whom you know that they possess sufficient intelligence and self-respect not to advance absurdities; to appeal to reason and not to authority, and to listen to reason and yield to it; and, finally, to cherish truth, to be willing to accept reason even from an opponent, and to be just enough to bear being proved to be in the wrong, should truth lie with him. From this it follows that scarcely one man in a hundred is worth your disputing with him.''

perfect words


\hfill

\end{sffamily}
\end{footnotesize}

